%------------------------------------------------------------------------------%

\usepackage{calculo_varias_variaveis-1}

\title{Curvas de Nível}

%------------------------------------------------------------------------------%
\begin{document}

\addtitle

%------------------------------------------------------------------------------%
\section{Curvas de Nível}

%------------------------------------------------------------------------------%
\begin{frame}{Curvas de Nível}
  \emph{Curva de nível}: conjunto de pontos no plano $xy$ onde
  \[
    f(x, y) = c
  \]
  para uma constante $c$
\end{frame}

%------------------------------------------------------------------------------%
\begin{frame}{Superfície de Nível}
  \emph{Superfície de nível}: conjunto de pontos onde
  \[
  f(x, y, z) = c
  \]
  para uma constante $c$
\end{frame}

\framedgraphic{Exemplo}{exemplo_curva_nivel_relevo}{}
\framedgraphic{Exemplo}{exemplo-curva-nivel-mapa}{}
\framedgraphic{Exemplo}{exemplo-contour-plot-1}{}
\framedgraphic{Exemplo}{exemplo-contour-plot-2}{}
\framedgraphic{Exemplo}{exemplo-superficie-nivel-esfera}{}
\framedgraphic{Exemplo}{exemplo-isosurface-pressao}{}
\framedgraphic{Exemplo}{exemplo_isofurface-cubo}{}
\framedgraphic{Exemplo}{exemplo-carro-cfd}{}
\framedgraphic{Exemplo}{exemplo-cfd-1}{}

%------------------------------------------------------------------------------%
\section{Exemplos}

% 01
%------------------------------------------------------------------------------%
\addtocounter{exemplo}{1}
\begin{frame}{Exemplo \theexemplo}
  Descreva e esboce as curvas de nível da função \(f(x, y) = 100 - x^2 - y^2\)
\end{frame}

\begin{frame}{Exemplo \theexemplo}
  \onslide<+->{Para cada valor $c\leq 100$}

  \begin{align*}
    \onslide<+->{f(x, y) & = c \\}
    \onslide<+->{100 - x^2 - y^2 = c \\}
    \onslide<+->{- x^2 - y^2 = c - 100 \\}
    \onslide<+->{x^2 + y^2 = 100 - c}
  \end{align*}

  \onslide<+->{As curvas de nível são circunferências centradas nas origem de
    raio $\sqrt{100-c}$}
\end{frame}

% 02
%------------------------------------------------------------------------------%
\addtocounter{exemplo}{1}
\begin{frame}{Exemplo \theexemplo}
  Descreva as superfícies de nível da função \(f(x, y, z) = \sqrt{x^2 + y^2 + z^2}\)
\end{frame}

\begin{frame}{Exemplo \theexemplo}
  \onslide<+->{Para cada valor $c \geq 0$}

  \begin{align*}
    \onslide<+->{f(x, y, z) & = c \\}
    \onslide<+->{\sqrt{x^2 + y^2 + z^2} & = c \\}
    \onslide<+->{x^2 + y^2 + z^2 & = c^2}
  \end{align*}

  \onslide<+->{As superfícies de nível são esferas centradas na origem de raio $c$}
\end{frame}

% 03
%------------------------------------------------------------------------------%
\addtocounter{exemplo}{1}
\begin{frame}{Exemplo \theexemplo}
  Descreva e esboce as curvas de nível da função \(f(x, y) = xy\)
\end{frame}

\begin{frame}{Exemplo \theexemplo}
  \onslide<+->{Para cada valor $c \neq 0$}

  \begin{align*}
    \onslide<+->{f(x, y) & = c \\}
    \onslide<+->{xy & = c \\}
    \onslide<+->{y & = \frac{c}{x}}
  \end{align*}

  \onslide<+->{Para $c=0$, ou $x=0$ ou $y=0$}
\end{frame}

% 04
%------------------------------------------------------------------------------%
\addtocounter{exemplo}{1}
\begin{frame}{Exemplo \theexemplo}
  Descreva e esboce as curvas de nível da função
  \(f(x, y) = \ln\left(9-x^2-y^2\right)\)
\end{frame}

\begin{frame}{Exemplo \theexemplo}
  \onslide<+->{Domínio de $f$}
  \begin{align*}
    \onslide<+->{9-x^2-y^2 & > 0         \\}
    \onslide<+->{9         & > x^2 + y^2 \\}
    \onslide<+->{x^2 + y^2 & < 3^2}
  \end{align*}
  \onslide<+->{Disco aberto centrado na origem de raio 3}

  \onslide<+->{Imagem:}
  \onslide<+->{a função assume valores no intervalo \(\big(-\infty, \ln(9)\,\big]\)}
\end{frame}

\begin{frame}{Exemplo \theexemplo}
  \onslide<+->{Para cada valor $c \leq \ln(9)$}

  \begin{align*}
    \onslide<+->{f(x, y)                   & = c       \\}
    \onslide<+->{\ln\left(9-x^2-y^2\right) & = c       \\}
    \onslide<+->{9-x^2-y^2                 & = e^c     \\}
    \onslide<+->{-x^2-y^2                  & = e^c - 9 \\}
    \onslide<+->{x^2+y^2                   & = 9 - e^c}
  \end{align*}

  \onslide<+->{As curvas de nível são circunferências centrada nas origem de
    raio $\sqrt{9 - e^c}$}
\end{frame}

% 05
%------------------------------------------------------------------------------%
\addtocounter{exemplo}{1}
\begin{frame}{Exemplo \theexemplo}
  Encontre a equação e esboce a curva de nível da função
  \[
    f(x, y) = \frac{2y-x}{x+y+1}
  \]
  que passa pelo ponto \((-1, 1)\)
\end{frame}

\begin{frame}{Exemplo \theexemplo}
  \begin{align*}
    \onslide<+->{c & = f(-1, 1)                              \\[2mm]}
    \onslide<+->{c & = \frac{2y-x}{x+y+1} \evalat{(-1, 1)}{} \\[2mm]}
    \onslide<+->{c & = \frac{2\times 1-(-1)}{(-1)+1+1}       \\[2mm]}
    \onslide<+->{c & = \frac{2+1}{1}                         \\[2mm]}
    \onslide<+->{c & = 3}
  \end{align*}
\end{frame}

\begin{frame}{Exemplo \theexemplo}
  \begin{align*}
    \onslide<+->{f(x, y)            & = 3        \\[2mm]}
    \onslide<+->{\frac{2y-x}{x+y+1} & = 3        \\[2mm]}
    \onslide<+->{2y-x               & = 3(x+y+1) \\[2mm]}
    \onslide<+->{2y-x               & = 3x+3y+3  \\[2mm]}
    \onslide<+->{2y-3y              & = 3x+x+3   \\[2mm]}
    \onslide<+->{-y                 & = 4x+3     \\[2mm]}
    \onslide<+->{y                  & = -4x -3}
  \end{align*}
\end{frame}

% 06
%------------------------------------------------------------------------------%
\addtocounter{exemplo}{1}
\begin{frame}{Exemplo \theexemplo}
  Encontre a equação e esboce a curva de nível da função
  \[
    f(x, y) = \sqrt{x^2-1}
  \]
  que passa pelo ponto \((1, 0)\)
\end{frame}

\begin{frame}{Exemplo \theexemplo}
Domínio
\begin{align*}
  \onslide<+->{x^2 -1  & \geq 0  \\[2mm]}
  \onslide<+->{x^2     & \geq 1  \\[2mm]}
  \onslide<+->{\abs{x} & \geq 1  \\[2mm]}
  \onslide<+->{x       & \leq -1 \quad\text{ou}\quad x \geq 1}
\end{align*}
\end{frame}

\begin{frame}{Exemplo \theexemplo}
  \begin{align*}
    \onslide<+->{c & = f(1, 0)                           \\[2mm]}
    \onslide<+->{c & = \sqrt{x^2-1} \; \evalat{(1, 0)}{} \\[2mm]}
    \onslide<+->{c & = \sqrt{1-1}                        \\[2mm]}
    \onslide<+->{c & = 0}
  \end{align*}
\end{frame}

\begin{frame}{Exemplo \theexemplo}
  \begin{align*}
    \onslide<+->{f(x, y)      & = 0 \\[2mm]}
    \onslide<+->{\sqrt{x^2-1} & = 0 \\[2mm]}
    \onslide<+->{x^2-1        & = 0 \\[2mm]}
    \onslide<+->{x^2          & = 1 \\[2mm]}
    \onslide<+->{\abs{x}      & = 1 \\[2mm]}
    \onslide<+->{x            & = \pm 1}
  \end{align*}
\end{frame}

% 07
%------------------------------------------------------------------------------%
\addtocounter{exemplo}{1}
\begin{frame}{Exemplo \theexemplo}
  Encontre uma equação para a superfície de nível da função
  \[
    f(x, y, z) = \frac{x - y + z}{2x + y -z}
  \]
  que passa pelo ponto \((1, 0, -2)\)
\end{frame}

\begin{frame}{Exemplo \theexemplo}
  Domínio
  \begin{align*}
    \onslide<+->{2x + y -z  & \neq 0  \\[2mm]}
    \onslide<+->{2x + y     & \neq z  \\[2mm]}
    \onslide<+->{z          & \neq 2x + y}
  \end{align*}
\end{frame}

\begin{frame}{Exemplo \theexemplo}
  \begin{align*}
    \onslide<+->{c & = f(1, 0, -2)                                          \\[2mm]}
    \onslide<+->{c & = \frac{x - y + z}{2x + y -z} \; \evalat{(1, 0, -2)}{} \\[2mm]}
    \onslide<+->{c & = \frac{1 - 0 -2}{2\times 1 + 0 -(-2)}                 \\[2mm]}
    \onslide<+->{c & = \frac{-1}{4}}
  \end{align*}
\end{frame}

\begin{frame}{Exemplo \theexemplo}
  \begin{align*}
    \onslide<+->{f(x, y, z)      & = -\frac{1}{4} \\[2mm]}
    \onslide<+->{\frac{x - y + z}{2x + y -z} & = -\frac{1}{4} \\[2mm]}
    \onslide<+->{-4(x - y + z)        & = 2x + y -z \\[2mm]}
    \onslide<+->{-4x +4y -4z        & = 2x + y -z \\[2mm]}
    \onslide<+->{-4z + z        & = 2x + y +4x -4y\\[2mm]}
    \onslide<+->{-3z        & = 6x -3y\\[2mm]}
    \onslide<+->{z        & = -2x + y}
  \end{align*}
\end{frame}

% 07
%------------------------------------------------------------------------------%
\addtocounter{exemplo}{1}
\begin{frame}{Exemplo \theexemplo}
Encontre a equação e esboce para a superfície de nível da função
\[
f(x, y, z) = \frac{x - y + z}{2x + y -z}
\]
que passa pelo ponto \((1, 0, -2)\)
\end{frame}

\begin{frame}{Exemplo \theexemplo}
Domínio
\begin{align*}
  \onslide<+->{2x + y -z  & \neq 0  \\[2mm]}
  \onslide<+->{2x + y     & \neq z  \\[2mm]}
  \onslide<+->{z          & \neq 2x + y}
\end{align*}
\end{frame}

\begin{frame}{Exemplo \theexemplo}
\begin{align*}
  \onslide<+->{c & = f(1, 0, -2)                                          \\[2mm]}
  \onslide<+->{c & = \frac{x - y + z}{2x + y -z} \; \evalat{(1, 0, -2)}{} \\[2mm]}
  \onslide<+->{c & = \frac{1 - 0 -2}{2\times 1 + 0 -(-2)}                 \\[2mm]}
  \onslide<+->{c & = \frac{-1}{4}}
\end{align*}
\end{frame}

\begin{frame}{Exemplo \theexemplo}
\begin{align*}
  \onslide<+->{f(x, y, z)                  & = -\frac{1}{4} \\[2mm]}
  \onslide<+->{\frac{x - y + z}{2x + y -z} & = -\frac{1}{4} \\[2mm]}
  \onslide<+->{-4(x - y + z)               & = 2x + y -z \\[2mm]}
  \onslide<+->{-4x +4y -4z                 & = 2x + y -z \\[2mm]}
  \onslide<+->{-4z + z                     & = 2x + y +4x -4y\\[2mm]}
  \onslide<+->{-3z                         & = 6x -3y\\[2mm]}
  \onslide<+->{z                           & = -2x + y}
\end{align*}
\end{frame}

%% 08
%%------------------------------------------------------------------------------%
%\addtocounter{exemplo}{1}
%\begin{frame}{Exemplo \theexemplo}
%  Encontre a equação e esboce para a superfície de nível da função
%  \[
%    f(x, y, z) = \sqrt{x-y} - \ln(z)
%  \]
%  que passa pelo ponto \((3, -1, 1)\)
%\end{frame}
%
%\begin{frame}{Exemplo \theexemplo}
%  \onslide<+->{Domínio}
%  \begin{align*}
%    \onslide<+->{x-y & \geq 0  \\[2mm]}
%    \onslide<+->{x   & \geq y}
%  \end{align*}
%  \onslide<+->{e
%  \[z > 0\]}
%
%\end{frame}
%
%\begin{frame}{Exemplo \theexemplo}
%  \begin{align*}
%    \onslide<+->{c & = f(3, -1, 1)                                  \\[2mm]}
%    \onslide<+->{c & = \sqrt{x-y} - \ln(z) \; \evalat{(3, -1, 1)}{} \\[2mm]}
%    \onslide<+->{c & = \sqrt{3-(-1)} - \ln(1)                       \\[2mm]}
%    \onslide<+->{c & = \sqrt{4}                                     \\[2mm]}
%    \onslide<+->{c & = 2}
%  \end{align*}
%\end{frame}
%
%\begin{frame}{Exemplo \theexemplo}
%  \begin{align*}
%    \onslide<+->{f(x, y, z)          & = 2 \\[2mm]}
%    \onslide<+->{\sqrt{x-y} - \ln(z) & = 2 \\[2mm]}
%    \onslide<+->{\sqrt{x-y}          & = 2 + \ln(z) \\[2mm]}
%    \onslide<+->{-4x +4y -4z         & = 2x + y -z \\[2mm]}
%    \onslide<+->{-4z + z        & = 2x + y +4x -4y\\[2mm]}
%    \onslide<+->{-3z        & = 6x -3y\\[2mm]}
%    \onslide<+->{z        & = -2x + y}
%  \end{align*}
%\end{frame}

%------------------------------------------------------------------------------%
\section{Lista Mínima}

%------------------------------------------------------------------------------%
\begin{frame}{Lista Mínima}
  Cálculo Vol. 2 do Thomas 12\textsuperscript{a} ed. -- Seção 14.1
  \begin{enumerate}
    \item Estudar o texto da seção
    \item Resolver os exercícios: 17-23, 31-36
  \end{enumerate}

  \vfill
  Atenção: A prova é baseada no livro, não nas apresentações
\end{frame}

%------------------------------------------------------------------------------%
\end{document}
%------------------------------------------------------------------------------%
